\section{引言}

大语言模型的数据来源往往横跨新闻、论坛、代码仓库和个人网站。模型参数是否记忆训练内容已经受到大量研究关注，但训练流程的第一个可学习组件——Tokenizer——长期被当作无敏感性的预处理器。BPE Tokenizer 将高频相邻符号反复合并，最终词表与合并序列因而是训练语料频率的压缩投影。Tokenizer 体积小、推理接口简单，常与闭源模型一起公开；攻击者不必承担语言模型前向计算成本，就能比较候选网站在目标 Tokenizer 下的压缩、词表和合并特征。近期工作已从 BPE 恢复语料混合比例\cite{hayase2024mixture}，并系统提出五类 Tokenizer 成员推断攻击\cite{tong2026tokenizermia}。

同态加密与差分隐私分别对应这一问题的两个层面。同态加密隐藏上传和聚合过程中的单客户端计数，但正确解密保持最终 Tokenizer 的输出分布；差分隐私通过随机化约束发布物对单个网站的依赖。本文因此把训练期机密性与发布期成员隐私拆成两个安全目标，并分别采用 Paillier 加法同态聚合与站点级差分隐私实现。

选择“站点”而非“文档”为隐私单位来自攻击问题本身：攻击者判断的是某个完整网站是否参与 Tokenizer 训练。网站级邻接关系直接刻画完整网站的加入或删除，避免 document-level DP 经群体隐私换算时随文档数放大预算。本文在每轮对一个网站的候选词对频率向量施加 L1 上界，使单轮 add/remove 站点敏感度可控，再采用整数双边几何机制。为减少逐次 BPE 带来的隐私组合与密码学交互，协议每轮选择最多 $b$ 个互不冲突合并。

本地差分隐私要求各站点在上传前独立加入噪声，其优势是不依赖可信聚合方，但多客户端噪声叠加可能降低高维词对统计的可用性。本文采用另一条协同防御路径：站点先在公共候选域上执行范数截断，再通过 Paillier 密文聚合隐藏原始计数，最后仅对聚合结果加入一次站点级噪声。该设计以非共谋双服务器为附加信任条件，换取更集中的噪声控制。Local-DP-BPE 作为对照，用于观察两种噪声位置对成员风险和下游效用的影响；由于两者采用不同隐私预算，该比较反映配置层面的隐私--效用差异，而非同预算严格对照。

本文的主要工作包括以下四个方面。第一，在固定版本的 C4 网站语料及互斥数据划分上，复现压缩率、词表重叠、频率估计、合并相似度和朴素贝叶斯五类分词器成员推断攻击，并报告 ROC AUC、低误报率性能和网站级 bootstrap 区间。第二，提出 SA-DP-BPE 协议，将公共候选、站点级截断、离散噪声和批量兼容合并统一到双服务器安全聚合框架中。第三，在 DCRA 和非共谋假设下分析 Paillier 聚合的计算机密性，并利用站点级敏感度与基本顺序组合给出发布机制的隐私保证。第四，通过 C4 分词效用、AG News 下游分类以及真实 2048-bit Paillier 开销实验，评估方法的隐私、效用和计算成本。
